\section{Executable tests and measurable predictive quotients}\label{app:quotient}
We give the elementary finite-closure construction underlying the representation in Section~\ref{sec:experiments}. Fix one prepared model. Let $\mathcal V$ be the linear space generated by bounded executable future-test functions on the physical state, including the constant function, the audit moments, and all legal one-step report pullbacks. A test is a specified causal continuation with a bounded final report, not an arbitrarily attached coordinate. Suppose $\mathcal V$ is finite-dimensional and admits the positive instrument representation \eqref{eq:raw-instrument}. Choose a basis $1,f_1,\ldots,f_d$ and evaluate it on prepared posterior measures.

\begin{proposition}[Realization and minimality]\label{prop:quotient}
Two prepared posterior measures give the same value to every test in $\mathcal V$ if and only if their basis evaluation vectors agree. Finite future report-and-action laws and the protected audit descend to this evaluation. On positive evidence, its update is \eqref{eq:update}. The history-to-evaluation map is Borel and gives a standard Borel realization in the compact base $D$. Any Borel statistic of a standard Borel history that determines all these future tests factors this evaluation through a Borel map.

This does not assert that the image of every Borel history map is a Borel subset of $D$. When only an analytic image is established, $D$ is the declared Borel realization and probability laws are pushforwards into $D$; unused points are not declared acquired. In the two raw realizations the post-report evaluation images are compact and therefore Borel.
\end{proposition}
\begin{proof}
Equality on a finite basis is equivalent to equality on its span. Pullback closure means that the probability of a report event and the unnormalized numerator of every subsequent test are linear evaluations of the incoming vector. Dividing by its positive report likelihood gives \eqref{eq:update}; zero likelihood has zero positive numerator by strict positivity of $e$ on the cone. Iterating gives every finite continuation law. Bounded monotone limits of such tests remain affine when they belong to the declared closure: evaluation on a fixed finite spanning family of prepared states determines the coefficients, so pointwise bounded convergence determines their limit. Alternatively, one may retain just the explicitly declared finite continuations; no closure under an unproved infinite-horizon operation is needed for the main theorem.

The posterior evaluation is Borel by standard Borel conditional kernels, or directly by finite matrix products and normalization on the finite-report history spaces, pieced together over their countable union. If another statistic $r(h)$ determines all tests, each evaluation coordinate is constant on its fibers. For completeness, the Borel factor statement follows by separating the two disjoint analytic images under $r$ of the Borel sets where a coordinate is, respectively, at most or greater than a fixed rational number. This shows that each coordinate is measurable with respect to $\sigma(r)$. The measurable factorization lemma for a real-valued function, applied coordinatewise, supplies a Borel factor; the resulting factor is defined on the entire standard Borel codomain of $r$. Its inverse image of the closed set $D$ is Borel and contains every attained statistic; replacing its values on the complement of that Borel set by a fixed point of $D$ makes it $D$-valued. No modification on a potentially non-Borel image complement is used. Thus the evaluation is the minimal information required by these tests, in this precise factorization sense.

The image caveat follows because Borel images of standard Borel spaces are in general only analytic. The general construction does not require the image to be the state space of the machine. In Theorem~\ref{thm:classical}, the map $x\mapsto\phi(x)$ is continuous and injective on compact $K$, and the audit probabilities separate the full convex base. In Theorem~\ref{thm:quantum}, the density matrices form a compact set and the informationally complete audit separates them. Their acquired rank orbits are also compact. These give the stronger image assertion for the realizations.
\end{proof}

For a family of unknown models, this construction is made within each model's finite predictive representation. A posterior on an unknown continuous model parameter can require additional, even infinite-dimensional, coordinates. Theorem~\ref{thm:main}(i) does not pretend otherwise: it couples the model-indexed certificates by a single observer kernel. It never gives that kernel access to a model-specific posterior. The geometric upper is stronger in a different respect: the current response $f_a(y)$ and the acceptance rule are common to the whole family.

\begin{example}[Separate feasibility does not give a common program]\label{ex:model-routing}
Prepare a fair bit $U$. In model $0$ the report is $Y=U$; in model $1$ it is $Y=1-U$. The protected audit is $U$. Both models have the same fair report law, and each separately has a perfect one-call decoder. No common decoder can be perfect: under the uniform prior on the two model labels, $U$ and $Y$ are independent fair bits. For squared $0/1$ prediction the common minimax risk is $1/4$, attained by reporting $1/2$. Maximizing the routing contribution separately within each model would incorrectly certify zero. This is why model contributions in \eqref{eq:joint-support} are added before its maximum, and why terminal readouts remain common.
\end{example}
